% from ega4-17-fr.tex
\oldpage[IV-4]{72}

\subsection{Propriétés de descente, de passage à la limite et de constructibilité}
\EGAReaderTarget{ega4.c04.subsec.17-7}{subsection:IV.17.7}
\label{IV.17.7}

\begin{proposition}[17.7.1]
\label{IV.17.7.1}
Soient $f:X\to Y$ un morphisme localement de présentation finie, $g:Y'\to Y$ un morphisme, $X'=X\times_Y Y'$, $f'=f_{(Y')}:X'\to Y'$ et $g':X'\to X$ les projections canoniques. Soit $x'$ un point de $X'$ et posons $x=g'(x')$, $y'=f'(x')$.
\begin{enumerate}
  \item[(i)] Si $f'$ est non ramifié au point $x'$, alors $f$ est non ramifié au point $x$.
  \item[(ii)] Supposons de plus que $g$ soit plat au point $y'$. Alors, si $f'$ est lisse (resp. étale) au point $x'$, $f$ est lisse (resp. étale) au point $x$.
\end{enumerate}
\end{proposition}

Posons $y=f(x)=g(y')$, de sorte que l'on a $f'^{-1}(y')=f^{-1}(y)\otimes_{k(y)}k(y')$; notons que $f'$ est localement de présentation finie.

(i) Comme la propriété pour un morphisme (localement de présentation finie) d'être non ramifié en un point ne fait intervenir que la fibre du morphisme en ce point (17.4.1, $d$)), on peut se borner au cas où $Y$ et $Y'$ sont des spectres de corps. Mais alors il revient au même de dire que $f$ (resp. $f'$) est non ramifié en $x$ (resp. $x'$) ou qu'il est étale en ce point (17.6.1, $c$)), donc (i) est une conséquence de (ii).

(ii) Comme $f'$ est plat au point $x'$ (17.5.1), l'hypothèse que $g$ est plat au point $y'$ entraîne que $f$ est plat au point $x$, car la projection $g':X'\to X$ est un morphisme plat au point $x'$, et $f\circ g'=g\circ f'$ est un morphisme plat au point $X'$, d'où la conclusion (2.2.11, (iv)). Le fait que $f$ soit lisse (resp. étale) au point $x$ ne fait plus alors intervenir que la fibre $f^{-1}(y)$ (17.5.1 et 17.6.1), et on est donc encore ramené au cas où $Y=\Spec(k)$ et $Y'=\Spec(k')$ sont des spectres de corps. Dire que $f'$ est lisse au point $x'$ signifie alors (17.5.1) que $X'$ est un $k'$-préschéma géométriquement régulier au point $x'$, et cela entraîne (6.7.8) que $X$ est un $k$-préschéma géométriquement régulier au point $x$, donc que $f$ est lisse au point $x$. Supposons de plus que $f'$ soit étale au point $x'$, de sorte que $x'$ est isolé dans $f'^{-1}(y')$ (17.6.1); comme la projection $f'^{-1}(y')\to f^{-1}(y)$ est un morphisme ouvert (2.4.10), $x$ est isolé dans $f^{-1}(y)$; comme on sait déjà que $f$ est lisse au point $x$, il est étale en ce point (17.6.1).

\begin{corollary}[17.7.2]
\label{IV.17.7.2}
\begin{enumerate}
  \item[(i)] Avec les notations de (17.7.1), soit $U$ (resp. $U'$) l'ensemble des points de $X$ (resp. $X'$) où $f$ (resp. $f'$) est non ramifié; alors on a $U'=g'^{-1}(U)$.
  \item[(ii)] Supposons de plus que $g$ soit plat, et soit $V$ (resp. $V'$) l'ensemble des points de $X$ où $f$ (resp. $f'$) est lisse; alors $V'=g'^{-1}(V)$.
\end{enumerate}
\end{corollary}

\begin{corollary}[17.7.3]
\label{IV.17.7.3}
\begin{enumerate}
  \item[(i)] Supposons $g$ surjectif; alors, pour que $f$ soit non ramifié, il faut et il suffit que $f'$ le soit.
  \item[(ii)] Soient $f:X\to Y$ un $S$-morphisme, $g:S'\to S$ un morphisme fidèlement plat. Supposons que $f$ soit localement de présentation finie, ou que $g$ soit quasi-compact. Alors, pour que $f$ soit lisse (resp. non ramifié, resp. étale), il faut et il suffit que $f'$ le soit.
\end{enumerate}
\end{corollary}

Dans (ii), le cas où $f$ est localement de présentation finie résulte de (17.7.1). Si $g$ est quasi-compact et $f'$ lisse (resp. non ramifié, resp. étale), $f'$ est localement de présentation finie par (2.7.1, (iv)) et on est ramené au premier cas.

\begin{proposition}[17.7.4]
\label{IV.17.7.4}
Soient $f:X\to Y$ un morphisme, $g:Y'\to Y$ un morphisme plat et localement de présentation finie, $X'=X\times_Y Y'$, $f'=f_{(Y')}:X'\to Y'$ et $g':X'\to X$ les

